\paragraph{归一化可见度、残余相干、经典化成本与解析腔室刚性}
在多重度剖面的 Fourier 能量下界、连续容量曲线的 layer-cake 对偶、实阶矩的 Mellin--Stieltjes 反演、全微态精确反演位率以及实参数无退化窗口已经建立之后，还可进一步把非零模可见度的算术障碍、残余相干的微分结构、经典化成本的可见度下界与解析读出的腔室刚性压缩为下列若干条结果。下文固定分辨率 \(m\)，并记
\[
M:=F_{m+2}=|X_m|.
\]
\[
G_m:=\ZZ/(M\ZZ),
\qquad
p_m(r):=\frac{d_m(r)}{2^m}
\qquad
(r\in G_m).
\]

\begin{theorem}[全模可见度同时湮灭的整除障碍]\label{thm:xi-time-part9m3-all-nonzero-visibility-extinction-divisibility-obstruction}
\leanverified{paper\_xi\_time\_part9m3\_all\_nonzero\_visibility\_extinction\_divisibility\_obstruction}
在该循环群上定义归一化字符振幅
\[
\mathcal A_m(k):=\sum_{r\in G_m}p_m(r)e^{-2\pi i kr/M}
=
\frac{\widehat d_m(k)}{2^m}
\qquad
(k\in G_m).
\]
则下列两条命题等价：
\[
\mathcal A_m(k)=0
\qquad
\text{对一切 }k\in G_m\setminus\{0\},
\]
\[
d_m(r)\equiv \frac{2^m}{M}
\qquad
\text{对一切 }r\in G_m.
\]
特别地，若 \(M\nmid 2^m\)，则必存在某个非零频率 \(k\in G_m\setminus\{0\}\) 使
\[
\mathcal A_m(k)\neq 0.
\]
进一步，当
\[
m\not\equiv 1\pmod 3
\]
时，\(F_{m+2}\) 为奇数，故必不整除 \(2^m\)；于是全部非零模同时湮灭在该分辨率族上被纯算术所禁止。
\end{theorem}
\begin{proof}
有限循环群上的 Fourier 反演给出
\[
p_m(r)=\frac1M\sum_{k\in G_m}\mathcal A_m(k)e^{2\pi i kr/M}.
\]
若对一切 \(k\neq 0\) 都有 \(\mathcal A_m(k)=0\)，则因
\[
\mathcal A_m(0)=\sum_{r\in G_m}p_m(r)=1,
\]
可得
\[
p_m(r)\equiv \frac1M,
\]
从而
\[
d_m(r)\equiv \frac{2^m}{M}.
\]
反向方向是平凡的：若 \(d_m\) 常值，则 \(p_m\) 为 \(G_m\) 上的均匀分布，其全部非平凡 Fourier 系数均为零。

若 \(M\nmid 2^m\)，则 \(\frac{2^m}{M}\notin\ZZ\)，从而上式常值条件不可能成立，因此至少有一个非零频率满足 \(\mathcal A_m(k)\neq 0\)。最后，Fibonacci 数列的奇偶性以周期 \(3\) 循环，故 \(F_n\) 为偶数当且仅当 \(3\mid n\)。于是当 \(m\not\equiv 1\pmod 3\) 时，\(M=F_{m+2}\) 为奇数；任何大于 \(1\) 的奇数都不可能整除 \(2^m\)，故结论成立。证毕。
\end{proof}

\begin{theorem}[归一化均方可见度的算术下界]\label{thm:xi-time-part9m3-mean-square-visibility-arithmetic-floor}
\leanverified{paper\_xi\_time\_part9m3\_mean\_square\_visibility\_arithmetic\_floor}
沿用定理 \ref{thm:xi-time-part9m3-all-nonzero-visibility-extinction-divisibility-obstruction} 的记号，并定义
\[
\Sigma_m^2
:=
\frac{1}{M-1}\sum_{k\in G_m\setminus\{0\}}|\mathcal A_m(k)|^2.
\]
将 \(2^m\) 除以 \(M\) 写成
\[
2^m=qM+s,
\qquad
0\le s<M.
\]
则有精确下界
\[
\Sigma_m^2\ge \frac{s(M-s)}{(M-1)2^{2m}}.
\]
并且等号成立，当且仅当多重度剖面达到唯一的最均匀整数配置：
\[
d_m(r)\in\{q,q+1\}
\qquad
(r\in G_m),
\]
且恰有 \(s\) 个剩余类满足 \(d_m(r)=q+1\)。
\end{theorem}
\begin{proof}
由定义
\[
|\mathcal A_m(k)|=\frac{|\widehat d_m(k)|}{2^m},
\]
故
\[
\Sigma_m^2
=
\frac{1}{(M-1)2^{2m}}\sum_{k\ne 0}|\widehat d_m(k)|^2.
\]
而结论 \ref{con:xi-time-part9l-fold-fourier-energy-arithmetic-lower-bound} 已给出精确能量下界
\[
\sum_{k\ne 0}|\widehat d_m(k)|^2\ge s(M-s),
\]
并且取等当且仅当 \(d_m\) 呈 \(q/q+1\) 两值配置，其中恰有 \(s\) 个剩余类取 \(q+1\)。代回上式即得所述不等式与取等刻画。证毕。
\end{proof}

\leanverified{paper\_xi\_time\_part9m3\_residual\_coherence\_differential\_calculus}
\begin{theorem}[残余相干曲线的微分结构]\label{thm:xi-time-part9m3-residual-coherence-differential-calculus}
\leanverified{paper\_xi\_time\_part9m3\_residual\_coherence\_differential\_calculus}
定义残余相干曲线
\[
\mathscr R_m(T):=
\frac{1}{2^m}\sum_{x\in X_m}(d_m(x)-T)_+^2
\qquad
(T\ge 0).
\]
则在几乎处处意义下有
\[
\mathscr R_m'(T)
=
-2^{1-m}\bigl(2^m-\mathcal C_m^{\mathrm{cont}}(T)\bigr),
\]
\[
\mathscr R_m''(T)=2^{1-m}N_m(T).
\]
从而
\[
\mathscr R_m(T)
=
2^{1-m}\int_T^\infty
\bigl(2^m-\mathcal C_m^{\mathrm{cont}}(s)\bigr)\,\dd s
\]
\[
=
2^{1-m}\int_T^\infty\int_s^\infty N_m(t)\,\dd t\,\dd s.
\]
此外，若记
\[
D_m^{\max}:=\max_{x\in X_m}d_m(x),
\qquad
\mathsf{Col}_m:=\sum_{r\in G_m}p_m(r)^2,
\]
则
\[
\mathscr R_m(0)=2^m\mathsf{Col}_m,
\qquad
\mathscr R_m(T)=0
\quad
\text{对一切 }T\ge D_m^{\max}.
\]
\end{theorem}
\begin{proof}
对任意固定整数 \(d\ge 0\)，函数
\[
f_d(T):=(d-T)_+^2
\]
满足
\[
f_d'(T)=-2(d-T)_+,
\qquad
f_d''(T)=2\,\mathbf 1_{\{T<d\}}
\]
在几乎处处意义下成立。逐项求和并除以 \(2^m\) 得
\[
\mathscr R_m'(T)
=
-2^{1-m}\sum_{x\in X_m}(d_m(x)-T)_+.
\]
另一方面，
\[
\sum_{x\in X_m}(d_m(x)-T)_+
=
\sum_{x\in X_m}\bigl(d_m(x)-\min(d_m(x),T)\bigr)
\]
\[
=
\sum_{x\in X_m}d_m(x)-\mathcal C_m^{\mathrm{cont}}(T)
=
2^m-\mathcal C_m^{\mathrm{cont}}(T),
\]
故第一条微分公式成立。

再由定理 \ref{thm:xi-capacity-tail-histogram-moment-complete-statistic}，
\[
\partial_T^+\mathcal C_m^{\mathrm{cont}}(T)=N_m(T)
\qquad
\text{a.e.},
\]
对上式再求导便得到
\[
\mathscr R_m''(T)=2^{1-m}N_m(T)
\qquad
\text{a.e.}
\]

由于 \(\mathscr R_m(T)=0\) 当 \(T\ge D_m^{\max}\) 时成立，故 \(\mathscr R_m(\infty)=0\)。把第一条公式从 \(T\) 积分到 \(+\infty\) 即得
\[
\mathscr R_m(T)
=
2^{1-m}\int_T^\infty
\bigl(2^m-\mathcal C_m^{\mathrm{cont}}(s)\bigr)\,\dd s.
\]
再由定理 \ref{thm:xi-time-part57b-capacity-moment-exact-mellin-stieltjes} 证明中已经出现的恒等式
\[
2^m-\mathcal C_m^{\mathrm{cont}}(s)
=
\int_s^\infty N_m(t)\,\dd t,
\]
代回即可得到双重积分表达。

最后，
\[
\mathscr R_m(0)
=
\frac{1}{2^m}\sum_{x\in X_m}d_m(x)^2
=
2^m\mathsf{Col}_m,
\]
而 \(T\ge D_m^{\max}\) 时每一项 \((d_m(x)-T)_+\) 都为零，因此 \(\mathscr R_m(T)=0\)。证毕。
\end{proof}

\begin{theorem}[容量、尾谱与残余相干的完全层析]\label{thm:xi-time-part9m3-capacity-residual-coherence-complete-tomography}
固定 \(m\)。记直方图
\[
h_m(n):=\#\{x\in X_m:\ d_m(x)=n\}
\qquad
(n\in\ZZ_{\ge 1}).
\]
则下列四类对象彼此一一决定：
\[
\{d_m(x):x\in X_m\},
\qquad
N_m,
\qquad
\mathcal C_m^{\mathrm{cont}},
\qquad
\mathscr R_m.
\]
更具体地，
\[
h_m(n)=N_m(n)-N_m(n+1),
\]
\[
N_m(T)=\partial_T^+\mathcal C_m^{\mathrm{cont}}(T)
\qquad
\text{a.e.},
\]
\[
N_m(T)=2^{m-1}\mathscr R_m''(T)
\qquad
\text{a.e.}
\]
因此，仅由 \(\mathcal C_m^{\mathrm{cont}}\) 或 \(\mathscr R_m\) 中任一条曲线，便可在固定分辨率 \(m\) 上唯一恢复全部纤维多重度多重集。进一步，对任意 \(q>0\)，都有
\[
S_q(m)=q\int_0^\infty t^{q-1}N_m(t)\,\dd t,
\]
故全部矩谱亦随之唯一恢复。
\leanverified{paper\_xi\_time\_part9m3\_capacity\_residual\_coherence\_complete\_tomography}
\end{theorem}
\begin{proof}
定理 \ref{thm:xi-capacity-tail-histogram-moment-complete-statistic} 已经给出
\[
\mathcal C_m^{\mathrm{cont}},
\qquad
N_m,
\qquad
h_m
\]
三者的完全互相恢复，并且同时给出
\[
S_q(m)=q\int_0^\infty t^{q-1}N_m(t)\,\dd t
\qquad(q>0).
\]
另一方面，定理 \ref{thm:xi-time-part9m3-residual-coherence-differential-calculus} 给出
\[
N_m(T)=2^{m-1}\mathscr R_m''(T)
\qquad
\text{a.e.}
\]
由于 \(N_m\) 是右连续、整数值且紧支撑的阶梯函数，其几乎处处取值已足以唯一确定整个函数；因此 \(\mathscr R_m\) 也唯一决定 \(N_m\)。

由
\[
h_m(n)=N_m(n)-N_m(n+1)
\]
可恢复直方图，而直方图与纤维多重度多重集显然互相等价。综上，四类对象彼此一一决定；再结合上式的 Mellin 表达，全部矩谱亦被唯一恢复。证毕。
\end{proof}

\begin{corollary}[最大纤维指数的精确经典化速率解释]\label{cor:xi-time-part9m3-exact-classicalization-maxfiber-rate}
\leanverified{paper\_xi\_time\_part9m3\_exact\_classicalization\_maxfiber\_rate}
将定理 \ref{thm:xi-full-microstate-exact-inversion-bitrate-threshold} 中的全微态精确反演预算解释为精确经典化成本。定义
\[
B_m^{\mathrm{cl}}
:=
\min\Bigl\{
B\in\ZZ_{\ge 0}:\ 2^B\ge d_m(x)\ \text{对一切 }x\in X_m
\Bigr\},
\]
并记
\[
D_m^{\max}:=\max_{x\in X_m}d_m(x).
\]
则
\[
B_m^{\mathrm{cl}}=\left\lceil \log_2 D_m^{\max}\right\rceil.
\]
若最大纤维指数
\[
\alpha_*:=\lim_{m\to\infty}\frac1m\log D_m^{\max}
\]
存在，则
\[
\lim_{m\to\infty}\frac{B_m^{\mathrm{cl}}}{m}
=
\frac{\alpha_*}{\log 2}.
\]
\end{corollary}
\begin{proof}
由定义，
\[
B_m^{\mathrm{cl}}
=
\min\{B\in\ZZ_{\ge 0}:2^B\ge D_m^{\max}\}
=
\left\lceil \log_2 D_m^{\max}\right\rceil.
\]
若 \(\alpha_*\) 存在，则
\[
\frac1m\log_2 D_m^{\max}
=
\frac{1}{m\log 2}\log D_m^{\max}
\longrightarrow
\frac{\alpha_*}{\log 2}.
\]
而上取整仅引入 \(O(1)\) 误差，故结论成立。证毕。
\end{proof}

\begin{theorem}[可见度与经典化成本的互补下界]\label{thm:xi-time-part9m3-visibility-classicalization-complementarity}
\leanverified{paper\_xi\_time\_part9m3\_visibility\_classicalization\_complementarity}
记 Rényi--\(2\) 有效宽度
\[
W_m^{(2)}:=\frac{1}{\mathsf{Col}_m}.
\]
则精确经典化成本满足
\[
B_m^{\mathrm{cl}}\ge m-\log_2 W_m^{(2)}.
\]
等价地，若用定理 \ref{thm:xi-time-part9m3-mean-square-visibility-arithmetic-floor} 中的均方可见度 \(\Sigma_m^2\) 表示，则
\[
B_m^{\mathrm{cl}}
\ge
m-\log_2 M+\log_2\bigl(1+(M-1)\Sigma_m^2\bigr).
\]
\end{theorem}
\begin{proof}
令
\[
p_{\max}:=\max_{r\in G_m}p_m(r).
\]
则
\[
\mathsf{Col}_m=\sum_{r\in G_m}p_m(r)^2
\le
p_{\max}\sum_{r\in G_m}p_m(r)
=
p_{\max}.
\]
因此
\[
D_m^{\max}=2^m p_{\max}\ge 2^m\mathsf{Col}_m=\frac{2^m}{W_m^{(2)}}.
\]
应用推论 \ref{cor:xi-time-part9m3-exact-classicalization-maxfiber-rate} 即得
\[
B_m^{\mathrm{cl}}
=
\left\lceil \log_2 D_m^{\max}\right\rceil
\ge
\log_2 D_m^{\max}
\ge
m-\log_2 W_m^{(2)}.
\]

另一方面，由 \(\mathcal A_m(0)=1\) 与 Parseval 恒等式，
\[
\sum_{k\in G_m}|\mathcal A_m(k)|^2
=
M\mathsf{Col}_m,
\]
故
\[
\Sigma_m^2
=
\frac{M\mathsf{Col}_m-1}{M-1}.
\]
解得
\[
\mathsf{Col}_m=\frac{1+(M-1)\Sigma_m^2}{M},
\]
代回第一式即得第二种写法。证毕。
\end{proof}

\begin{theorem}[容量亏损与残余相干的 Mellin 表示]\label{thm:xi-time-part9m3-capacity-defect-residual-coherence-mellin}
对任意实数 \(q>1\)，有
\[
S_q(m)
=
q(q-1)\int_0^\infty
T^{q-2}\bigl(2^m-\mathcal C_m^{\mathrm{cont}}(T)\bigr)\,\dd T.
\]
若进一步 \(q>2\)，则还有
\[
S_q(m)
=
q(q-1)(q-2)\,2^{m-1}
\int_0^\infty T^{q-3}\mathscr R_m(T)\,\dd T.
\]
\end{theorem}
\begin{proof}
第一条公式正是定理 \ref{thm:xi-time-part57b-capacity-moment-exact-mellin-stieltjes} 的第二条积分表达。

再由定理 \ref{thm:xi-time-part9m3-residual-coherence-differential-calculus}，
\[
2^m-\mathcal C_m^{\mathrm{cont}}(T)
=
-2^{m-1}\mathscr R_m'(T)
\qquad
\text{a.e.}
\]
故对 \(q>2\)，
\[
\begin{aligned}
S_q(m)
&=
q(q-1)\int_0^\infty
T^{q-2}\bigl(2^m-\mathcal C_m^{\mathrm{cont}}(T)\bigr)\,\dd T\\
&=
-q(q-1)\,2^{m-1}\int_0^\infty T^{q-2}\mathscr R_m'(T)\,\dd T.
\end{aligned}
\]
对最后一个积分分部积分。由于 \(\mathscr R_m\) 在 \([D_m^{\max},\infty)\) 上恒为零，且 \(q>2\) 时有
\[
T^{q-2}\mathscr R_m(T)\longrightarrow 0
\qquad
(T\downarrow 0),
\]
边界项全部消失，于是得到
\[
S_q(m)
=
q(q-1)(q-2)\,2^{m-1}
\int_0^\infty T^{q-3}\mathscr R_m(T)\,\dd T.
\]
证毕。
\end{proof}

\begin{proposition}[解析读出腔室内的折叠谱刚性]\label{prop:xi-time-part9m3-analytic-chamber-rigidity-sampling-readout}
\leanverified{paper\_xi\_time\_part9m3\_analytic\_chamber\_rigidity\_sampling\_readout}
固定 \(m\)，设 \(I\subset\RR\) 为区间。对每个微态 \(\omega\in\Omega_m\) 与每个标签 \(x\in X_m\)，设
\[
s_{\omega,x}:I\to\RR
\]
为实解析函数，并假定对每个 \(\theta\in I\) 都存在唯一最大者
\[
\Fold_{m,\theta}(\omega):=
\arg\max_{x\in X_m}s_{\omega,x}(\theta).
\]
定义对应纤维计数
\[
d_{m,\theta}(x):=
\#\{\omega\in\Omega_m:\ \Fold_{m,\theta}(\omega)=x\}.
\]
则存在离散集 \(E_m\subset I\)，使得对 \(I\setminus E_m\) 的每个连通分支 \(J\)，函数
\[
\theta\longmapsto d_{m,\theta}(x)
\]
对每个 \(x\in X_m\) 都在 \(J\) 上常值。因而
\[
D_{m,\theta}^{\max}:=\max_{x\in X_m}d_{m,\theta}(x),
\]
\[
\mathcal C_{m,\theta}^{\mathrm{cont}}(T):=
\sum_{x\in X_m}\min\bigl(d_{m,\theta}(x),T\bigr),
\]
\[
\mathscr R_{m,\theta}(T):=
\frac1{2^m}\sum_{x\in X_m}(d_{m,\theta}(x)-T)_+^2,
\]
\[
S_q(m;\theta):=\sum_{x\in X_m}d_{m,\theta}(x)^q
\qquad
(q>0)
\]
都在 \(J\) 上保持常值。

进一步，若参数 \(\theta\) 取自真实输入 \(40\) 状态核输出位势的实对数坐标，则注 \ref{rem:real-input-40-output-potential-degeneracy-phase} 给出统一窗口
\[
|\theta|<\theta_0,
\qquad
\theta_0=0.0533477827\ldots,
\]
其中不会跨越正实退化点。因而任何可归约为上述解析 \(\arg\max\) 读出的实参数变形，只有在击中离散交叉集 \(E_m\) 时才可能改变折叠谱与由其诱导的容量、残余相干及矩谱。
\end{proposition}
\begin{proof}
对每个三元组 \((\omega,x,y)\)（其中 \(x\neq y\)），考虑解析差值函数
\[
f_{\omega,x,y}(\theta):=s_{\omega,x}(\theta)-s_{\omega,y}(\theta).
\]
若 \(f_{\omega,x,y}\not\equiv 0\)，则其零点在 \(I\) 内离散。定义
\[
E_m:=
\bigcup_{\omega\in\Omega_m}
\bigcup_{\substack{x,y\in X_m\\ x\neq y\\ f_{\omega,x,y}\not\equiv 0}}
\{\theta\in I:\ f_{\omega,x,y}(\theta)=0\}.
\]
由于 \(\Omega_m\) 与 \(X_m\) 都有限，\(E_m\) 是有限个离散集之并，故仍为离散集。

取 \(I\setminus E_m\) 的任一连通分支 \(J\)。若对某个固定 \(\omega\) 而言，\(\Fold_{m,\theta}(\omega)\) 在 \(J\) 上发生变化，则存在 \(\theta_1<\theta_2\) 属于 \(J\) 与两个不同标签 \(x\neq y\)，使得
\[
\Fold_{m,\theta_1}(\omega)=x,
\qquad
\Fold_{m,\theta_2}(\omega)=y.
\]
由连续性，介值定理迫使某个内点 \(\theta_\ast\in J\) 上有
\[
s_{\omega,x}(\theta_\ast)=s_{\omega,y}(\theta_\ast),
\]
即 \(f_{\omega,x,y}(\theta_\ast)=0\)。若该差值函数不恒为零，则 \(\theta_\ast\in E_m\)，与 \(J\cap E_m=\varnothing\) 矛盾；若该差值函数恒为零，则 \(x\) 与 \(y\) 在整个 \(I\) 上始终并列，从而违背“唯一最大者”的假设。故对每个 \(\omega\)，标签 \(\Fold_{m,\theta}(\omega)\) 在 \(J\) 上恒定。

于是对每个 \(x\in X_m\)，计数 \(d_{m,\theta}(x)\) 在 \(J\) 上恒定。其余各量都只是 \(d_{m,\theta}(x)\) 的确定性函数，故亦在 \(J\) 上恒定。最后一段只是在真实输入 \(40\) 输出位势的实参数坐标中调用注 \ref{rem:real-input-40-output-potential-degeneracy-phase} 所给出的无退化窗口。证毕。
\end{proof}

\endinput
